%=============================================================
\section{Calibration Mechanism: Implementation Details and Proofs}
\label{app:calibration_details}
%=============================================================

This appendix supplies the technical details deferred from
\S\ref{sec:calibration}. We provide necessary justifications, document every adaptive parameter and how
it is initialized at cold start, prove the recovery property of the
credibility update, justify the latency claim, characterize the
additional memory and token overhead, and give implementation notes
for both the embedding and graph instantiations of \system.
\subsection{Mapping Between Calibration Quantities and Utility Components}
\label{app:utility-mapping}

This subsection makes precise the claim, illustrated visually in
Figure~\ref{fig:method}, that every component of the contributor utility
$u_c = \Pr[C{=}1]\cdot I - D$ and the auditor utility
$u_i = \mathrm{Align}(\Delta_t,M_t) - c_i$ defined in
Section~\ref{sec:problem} is controlled, at the calibrated equilibrium of
$\widetilde\Gamma$, by exactly one of the three calibration quantities in
\system: the detection signal $\rho_{\text{detect}}$, the alignment
signal $\rho_{\text{align}}$, and the credibility weight $w_i^{(t)}$
(into which we fold the commit gate $\mathbf{1}[\rho \geq \rho^*]$, since
the gate's structural authority comes from $w_i^{(t)}$ entering $\rho$
multiplicatively). Table~\ref{tab:utility-mapping} summarizes the mapping;
the rest of the subsection states and justifies each row. We address
mechanisms only insofar as they bear on equilibrium utility control;
implementation details are deferred to the rest of Appendix \ref{app:calibration_details}.

\begin{table}[h]
\centering
\small
\caption{Per-component mapping between the three calibration quantities
of \system and the utility components they control at the calibrated
equilibrium. Each row identifies the component, the calibration quantity
that controls it, and the property that does the control.}
\label{tab:utility-mapping}
\begin{tabular}{@{}llll@{}}
\toprule
Calibration quantity & Side & Utility component & Controlling property \\
\midrule
\multirow{2}{*}{$\rho_{\text{detect}}$}
 & Contributor & $D(\Delta_t,M_t)$ & detection regulation \\
 & Contributor & $I(\Delta_t,M_t)$ & influence ceiling \\
\midrule
\multirow{2}{*}{$\rho_{\text{align}}$}
 & Auditor     & $\mathrm{Align}(\Delta_t,M_t)$ & sample-time regulation \\
 & Auditor     & $c_i$ & probe re-use and damping \\
\midrule
\multirow{2}{*}{$w_i^{(t)}$ (with commit gate)}
 & Contributor & $\Pr[C{=}1]$ & vote-independent regulation \\
 & Auditor     & $c_i$ &  damping \\
\bottomrule
\end{tabular}
\end{table}

\subsubsection{The detection signal \texorpdfstring{$\rho_{\text{detect}}$}{rho\_detect} controls \texorpdfstring{$D$}{D} and \texorpdfstring{$I$}{I}}

The detection signal scores how the proposed $\Delta_t$ sits in the
geometry (embedding case) or topology (graph case) of $M_t$. Both
contributor utility components on the left of $u_c$, namely $D$ and
$I$, are controlled by this single quantity at equilibrium.

\textbf{Control of $D(\Delta_t,M_t)$ via detection regulation.}
In the uncalibrated game $\Gamma$, $D$ is conceptual: detection has no
operational meaning unless an external mechanism makes inconsistency
observable. \system supplies that mechanism by setting
$D := 1 - \rho$, where $\rho_{\text{detect}}$ is the dominant input
(low alignment alone cannot redeem a high-density delta because
$\rho = \sqrt{\rho_{\text{detect}}\cdot\rho_{\text{align}}}$).

\textit{Justification.}
By construction $\rho \in [0,1]$, hence $D \in [0,1]$ as required by
Eq.~(\ref{eq:utility-contributor}). The substitution $D = 1-\rho$
is the natural one in the sense that the correlation coefficient
between $D$ and $1-\rho$ is exactly $\kappa = 1$, achieving the
manipulation ceiling derived in Appendix~\ref{app:failure_modes}
(Channel~1). In particular, $\rho_{\text{detect}}$ becomes informative
as a detection signal precisely when $\Delta_t$ disturbs the local
density (embedding) or violates local/path consistency (graph) -- the
two manipulations the contributor's $I$ relies on. There is therefore
no slack at equilibrium in which a high-$I$ delta scores low on
$D$.

\textbf{Control of $I(\Delta_t,M_t)$ via influence ceiling.}
The influence term $I$ measures how much $\Delta_t$ shifts downstream
retrieval. A manipulative contributor maximizes $I$ by placing
$\Delta_t$ in regions of $M_t$ that dominate retrieval, i.e., regions of
high local embedding density, or graph edges that bypass existing
multi-hop paths. These are exactly the patterns that
$\rho_{\text{detect}}$ scores; the retrieval re-scaling described in
\S\ref{ssec:retrieval} then forces every committed entry's effective
influence to be capped by its $\rho_{\text{detect}}$ score.

\textit{Justification (embedding case).}
For an entry $e_k^\Delta$,
$\rho_{\text{detect}}(\Delta_t) = 1 - \frac{1}{K}\sum_k \hat r_k/(\bar r+\epsilon)$.
Let $I_k$ denote the marginal influence of $e_k^\Delta$ on top-$k$
retrieval against an arbitrary query distribution. Top-$k$ retrieval
returns the $k$ nearest neighbours, so $e_k^\Delta$ enters the result
for any query lying within a ball of radius proportional to its
neighbourhood density, giving $I_k = \Theta(\hat r_k)$ asymptotically
in $|M_t|$. The retrieval re-scaling
$\tilde e_k = \sqrt{\rho_k}\,e_k$ at read time then yields
\[
I_k^{\text{eff}} \;=\; \rho_k \cdot I_k \;\leq\; \rho_k \cdot \Theta(\hat r_k)
\]
which is jointly capped because $\rho_k$ is decreasing in $\hat r_k$
through $\rho_{\text{detect}}$. A delta cannot simultaneously score
high on $\rho_{\text{detect}}$ and exert large $I$.

\textit{Justification (graph case).}
For a new edge $e=(u,r,v)$, $I(e)$ is determined by how often $e$ lies
on multi-hop retrieval paths. The path-strength product
$\prod_{e\in\pi}\rho_e$ used at retrieval ensures that any path through
a low-$\rho_{\text{detect}}$ edge is automatically discounted, so
$I^{\text{eff}}(e) \leq \rho_{\text{detect}}(e)\cdot I_{\max}$. The
ceiling $D \leq 1/(1+\kappa)$ in
Eq.~(\ref{eq:integrity_gap_bound}) of
Appendix~\ref{app:gap_closing} then gives a closed-form upper bound on
equilibrium influence.

\subsubsection{The alignment signal \texorpdfstring{$\rho_{\text{align}}$}{rho\_align} controls \texorpdfstring{$\mathrm{Align}$}{Align} and \texorpdfstring{$c_i$}{c\_i}}

The alignment signal is computed from auditor probes (queries $q_i$ for
embedding memory, traversal paths $\pi_i$ for graph memory). Both
auditor utility terms are controlled by this single quantity at
equilibrium: the alignment signal in $u_i$ is regulated in
$\rho_{\text{align}}$ directly, and the effort cost $c_i$ collapses
because the probes that produce $\rho_{\text{align}}$ are reused from
the auditor's normal debate activity.

\textbf{Control of $\mathrm{Align}(\Delta_t,M_t)$ via sample-time non-manipulability.}
\system regulates the alignment signal as
$\mathrm{Align}(\Delta_t,M_t) := 2\rho - 1 \in [-1,+1]$
(Eq.~\ref{eq:utility-auditor}, Appendix~\ref{app:utility_rationale}),
where $\rho_{\text{align}}$ is the auditor-side input. We make
explicit why this regulating is unmanipulable by any auditor at
equilibrium.

\textit{Justification.}
Let auditor $a_i$ submit probe evidence $\xi_i$. The alignment signal
decomposes as
\[
\rho_{\text{align}}(\Delta_t)
\;=\; 1 - \frac{1}{|\mathcal{A}_t|}\sum_{j\in\mathcal{A}_t} w_j^{(t)}\cdot s(\xi_j;\Delta_t,M_t)
\]
where $s$ is the per-auditor disagreement (RBO distance for embeddings,
indicator of unreachability for graphs). Two structural properties
make $\rho_{\text{align}}$ robust to single-auditor manipulation.
First, $\xi_i$ is sampled uniformly from $Q_i$ (or fixed to the
auditor's current path), so $a_i$ cannot pre-commit to a probe that
flatters its preferred outcome; the sampling is performed by \system,
not $a_i$. Second, the multiplicative entry of $w_j^{(t)}$ caps any
single auditor's marginal influence:
\[
\Big|\frac{\partial \rho_{\text{align}}}{\partial \xi_i}\Big| \;\leq\; \frac{w_i^{(t)}}{|\mathcal{A}_t|}
\]
At equilibrium, honest auditors converge to $w_i^{(t)}\to 1$
(Proposition~\ref{prop:credibility_fp}), giving them full influence;
adversarial auditors converge toward $0$. The alignment term in $u_i$
therefore tracks the structural quantity $2\rho-1$ rather than any
self-reported judgment.

\textbf{Damping $c_i$ via probe re-use.}
The effort cost $c_i$ in $u_i$ models the resource cost an auditor
incurs by performing scrutiny. In $\Gamma$, scrutiny is a separate
costly action, which produces the free-riding equilibrium
(Proposition~\ref{prop:freeriding}). \system neutralizes this cost
because the probes used to compute $\rho_{\text{align}}$ are exactly
the queries and traversals the auditor already issues during debate.

\textit{Justification.}
Let $\mathcal{T}_i^{\text{debate}}$ be the set of LLM calls and tool
invocations $a_i$ performs as part of normal debate participation, and
let $\mathcal{T}_i^{\Phi}$ be the additional set induced by \system.
Every probe used by \system is sampled from $Q_i$ (which $a_i$
generates regardless of \system) or set to $\pi_i$ (the path $a_i$
already walked), so
\[
\mathcal{T}_i^{\Phi} \subseteq \mathcal{T}_i^{\text{debate}}
\]
The empirical effort cost (Appendix~\ref{app:utility_rationale})
\[
\widehat c_i^{(t)} \;=\; \frac{1}{W}\sum_{t'=t-W}^{t-1}\big[\text{cost}(\mathcal{T}_i^{\text{debate}} \cup \mathcal{T}_i^{\Phi}, t') - \text{cost}(\mathcal{T}_i^{\text{debate}}, t')\big] \;=\; 0
\]
in expectation, modulo $O(1)$ algorithmic overhead dominated by LLM
inference cost (Appendix~\ref{app:latency}). The marginal scrutiny cost
is therefore asymptotically zero at equilibrium, and the free-riding
incentive of Proposition~\ref{prop:freeriding} dissolves.

\subsubsection{The credibility weight \texorpdfstring{$w_i^{(t)}$}{w\_i} (with commit gate) controls \texorpdfstring{$\Pr[C{=}1]$}{Pr[C=1]} and \texorpdfstring{$c_i$}{c\_i}}

The credibility weight enters $\rho$ multiplicatively through
$\rho_{\text{align}}$ and gates commitment via
$\mathbf{1}[\rho \geq \rho^*]$. The commit gate is therefore not a
separate calibration quantity but the structural endpoint of $w_i^{(t)}$:
without credibility weighting, the gate would be vote-equivalent. We
group them and show that together they control the contributor's
commit-probability term and the auditor's collusion incentive at
equilibrium.

\textbf{Control of $\Pr[C{=}1]$ via vote-independent commitment.}
Under $\Gamma$, the commit probability is the expectation over auditor
strategies $\Pr[C(\{v_i\})=1\mid\Delta_t]$, which depends multilinearly
on the mixed strategies $\{\pi_i\}$ and is therefore manipulable by any
pivotal coalition (Propositions~\ref{prop:freeriding}--\ref{prop:adversarial_rejection}).
Under $\widetilde\Gamma$, \system replaces this term with the indicator
$\mathbf{1}[\rho \geq \rho^*]$. The structural authority of this
substitution comes from $w_i^{(t)}$: an adversarial auditor who tries
to push $\rho_{\text{align}}$ up or down sees its $w_i^{(t)}$ decay
multiplicatively, so its leverage on the commit decision vanishes at
equilibrium.

\textit{Justification.}
Fix any $\Delta_t$ and $M_t$, and let
$\sigma=(\sigma_c, \sigma_1, \dots, \sigma_{|\mathcal{A}_t|})$ be the
joint strategy profile. The commit indicator under $\widetilde\Gamma$
satisfies
\[
\frac{\partial}{\partial \sigma_i} \,\mathbf{1}[\rho(\Delta_t,M_t)\geq\rho^*] \;=\; 0
\qquad\forall\, a_i\in\mathcal{A}_t
\]
because $\rho$ is computed from $(\Delta_t,M_t)$ and from the sampled
probes, neither of which is a vote or self-report; the auditor's only
channel of influence is its probe distribution, weighted by
$w_i^{(t)}$. At the calibrated equilibrium, $w_i^{(t)} \to 1 - p_i$
where $p_i$ is the auditor's long-run invalid-probe rate
(Proposition~\ref{prop:credibility_fp}), so an adversarial auditor's
contribution to the commit decision is bounded by $1 - p_i$ and shrinks
to zero as it persists in attempting manipulation. Conditional on prior
debate actions, the commit decision is therefore a deterministic
function of $(\Delta_t,M_t)$ and an auditor's marginal vote does not
move it. This eliminates the pivot-probability term that drives the
three failure equilibria and replaces vote-based commitment with a
structural rule.

\textbf{Control of $c_i$ via multiplicative damping.}
The cost $c_i$ in $u_i$ models the resource effort an auditor incurs
when submitting a structurally invalid probe in support of a manipulative
$\Delta_t$ (collusion). The credibility weight $w_i^{(t)}$ damps the
auditor's lifetime alignment payoff multiplicatively whenever such an
invalid probe is detected, so collusion incurs both an immediate cost
and a long-run weight loss.

\textit{Justification.}
A single invalid probe at round $t$ multiplies the auditor's
subsequent credibility trajectory by a factor of $(1-\delta)$ relative
to honest play (Eq.~\ref{eq:cred_weight}). Consider a one-shot
deviation at round $t=0$ followed by honest behaviour. Let
$\bar A := \mathbb{E}[\mathrm{Align}(\Delta_t,M_t)]$ at the calibrated
equilibrium and $\bar w := \mathbb{E}[w_i^{(t)}]\to 1$ along the honest
trajectory. The discounted lifetime payoff is
\[
V_i \;=\; \sum_{t=0}^{\infty}\beta^t\!\left(\mathrm{Align}(\Delta_t,M_t)\cdot w_i^{(t)}\;-\;c_i\cdot\mathbf{1}[\text{collude at }t]\right)
\]
Computing the difference between honest and colluding play under a
one-shot deviation,
\[
V_i^{\text{honest}}-V_i^{\text{collude}} \;=\; c_i \;+\; \sum_{t=1}^{\infty}\beta^t\,\bar A\,\bar w\,\delta
\;=\; c_i \;+\; \frac{\beta\,\delta}{1-\beta}\,\bar A\,\bar w \;>\; 0
\]
for any $c_i\geq 0$, $\bar A>0$, and $\beta\in(0,1)$. The first term is
the immediate cost the auditor avoids by staying honest; the second is
the long-run weight loss that compounds geometrically because
credibility decay propagates through every future round's alignment
contribution. Both terms favour honesty, so collusion is strictly
dominated under $\widetilde\Gamma$ at the calibrated equilibrium without $c_i$ needing any specific value: even at $c_i=0$, the
multiplicative damping alone suffices.

%-------------------------------------------------------------
\subsection{Adaptive Parameters and Cold-Start Behavior}
\label{app:adaptive_params}
%-------------------------------------------------------------

\system uses several adaptive parameters in place of fixed
hyperparameters. Each is computed from observable round statistics,
so no value is set by the user. Table~\ref{tab:adaptive_params}
summarizes them.

\begin{table}[h]
\caption{Adaptive parameters used by \system. Each is computed from
observable round statistics rather than tuned.}
\centering
\small
\begin{tabular}{lll}
\toprule
\textbf{Symbol} & \textbf{Definition} & \textbf{Source} \\
\midrule
$\rho^*$ & commit gate & running median of last $W$ committed
$\rho$ scores \\
$\delta_t$ & credibility update rate & fraction of probes invalid in
round $t$ \\
$\eta_t$ & path-compatibility cutoff & running median of past
$\text{comp}(\cdot)$ values \\
$\alpha_t$ & local/path mixture weight & fraction of edges in
$\Delta_t$ with $|\mathcal{P}_{uv}| > 0$ \\
adaptive radius & embedding similarity cutoff & running median of
pairwise similarities in $\mathcal{M}_t$ \\
$L_{\max}$ & BFS depth bound & $\lceil \log |\mathcal{V}_t| \rceil$ \\
$W$ & history window for medians & $\max(30, \lceil \sqrt{|\mathcal{M}_t|}
\rceil)$ \\
\bottomrule
\end{tabular}
\label{tab:adaptive_params}
\end{table}

\textbf{Cold start for the commit gate $\rho^*$.}
\label{app:rho_star}
At round $t=0$ no committed scores exist, so the running median is
undefined. We initialize $\rho^* = 0.5$, the midpoint of the score
range. Every $\Delta_t$ is provisionally accepted while the history
buffer fills. Once the buffer reaches $W$ entries, $\rho^*$ becomes
the running median and \system enters its steady-state regime. This
warm-up affects only the first $W$ rounds, and the only effect is
that early commits may be lower-trust than steady-state commits.
Subsequent retrieval discounts these early entries through the
trust-weight mechanism (\S\ref{ssec:retrieval}), so any inaccuracy is
absorbed at read time rather than persisting in the memory.

\textbf{Cold start for the credibility rate $\delta_t$.}
At $t = 0$ no probes have been observed, so we initialize $\delta_0 =
0.1$, a conservative default that produces a slow update. From round
$t = 1$ onward, $\delta_t$ is the empirical fraction of invalid
probes observed in round $t$. The choice of $\delta_0$ affects only
the first round's weight movement and is forgotten by round $t = 2$,
so it is not a tunable parameter in any meaningful sense.

\textbf{Cold start for $\eta_t$ and $\alpha_t$.}
Both depend on observed graph statistics. At $t = 0$ we set $\eta_0 =
0$ (accept any compatibility score) and $\alpha_0 = 0.5$ (equal
weighting of local and path checks). They become data-driven from
round $1$ onward.

\textbf{Why these defaults are not tunable.}
A reviewer might object that $\rho^*_0 = 0.5$, $\delta_0 = 0.1$,
$\eta_0 = 0$, and $\alpha_0 = 0.5$ are themselves choices. We
emphasize that these defaults affect only a $O(W)$-round warm-up and
have no effect on steady-state behavior, since the running statistics
overwrite each initial value within $W$ rounds.

%-------------------------------------------------------------
\subsection{Credibility Recovery: Proof of Fixed Point}
\label{app:credibility}
%-------------------------------------------------------------

We prove the claim from \S\ref{ssec:arch_calib} that an auditor whose
long-run invalid-probe rate is $p$ has credibility weight converging
to $w^* = 1 - p$.

\begin{proposition}[Credibility fixed point]
\label{prop:credibility_fp}
Let $a_i$ produce invalid probes independently with probability $p
\in [0,1]$ at each round, and let $\delta_t \in (0, 1)$ for all $t$.
Then under Eq.~\ref{eq:cred_weight}, $\mathbb{E}[w_i^{(t)}] \to 1 - p$
as $t \to \infty$.
\end{proposition}

\begin{proof}
Taking expectations on both sides of Eq.~\ref{eq:cred_weight} and
treating $\delta_t$ as approximately constant at its long-run value
$\delta$ (justified because $\delta_t$ converges to its empirical
mean by the law of large numbers):
\begin{align}
\mathbb{E}[w_i^{(t+1)}]
&= p \cdot \mathbb{E}[w_i^{(t)}](1 - \delta) +
   (1-p) \left( \mathbb{E}[w_i^{(t)}] + \delta(1 - \mathbb{E}[w_i^{(t)}]) \right) \\
&= \mathbb{E}[w_i^{(t)}] \left( p(1-\delta) + (1-p)(1 - \delta) \right) +
   (1-p) \delta \\
&= \mathbb{E}[w_i^{(t)}] (1 - \delta) + (1-p) \delta.
\end{align}
This is a linear contraction with contraction factor $1 - \delta \in
(0, 1)$ and offset $(1-p) \delta$. Its unique fixed point is
\begin{equation}
w^* = \frac{(1-p) \delta}{1 - (1 - \delta)} = 1 - p,
\end{equation}
and convergence is geometric at rate $1 - \delta$. \qed
\end{proof}

\textbf{Recovery time.}
The half-life of weight movement is $\log 2 / \log(1/(1-\delta))
\approx \ln 2 / \delta$ rounds. With a typical $\delta \approx 0.1$,
an honest auditor that suffered a one-time slip recovers half the
distance to $w^* = 1$ in roughly $7$ rounds, and reaches within $5\%$
of $w^*$ in roughly $30$ rounds.

\textbf{Robustness to time-varying $\delta_t$.}
If $\delta_t$ is non-constant but bounded away from $0$ and $1$, the
update remains a contraction and the fixed point still tracks
$1 - p_t$, where $p_t$ is the auditor's invalid rate at round $t$.
This means a previously dishonest auditor who reforms can recover
full credibility, and a previously honest auditor who turns
adversarial loses credibility without instant collapse, both of
which are desirable robustness properties.

%-------------------------------------------------------------
\subsection{Latency Analysis}
\label{app:complexity}
%-------------------------------------------------------------

We show that calibration adds only marginal time on top of the
baseline retrieval cost in both architectures. The key reason is that
\system reuses computations and data structures already produced by
baseline retrieval, so it never pays a second time for the same work.

\textbf{Embedding case.}
Baseline top-$k$ retrieval costs $O(k \log N)$ per query using an
HNSW or IVF index. Calibration adds two passes whose cost can be
analyzed below the baseline:
\begin{itemize}
    \item \textbf{Detection pass.} For each of $K$ new entries we
    perform an ANN range query to count neighbors within the
    adaptive radius. Because $K \ll N$ in any incremental update
    (the delta is small relative to memory size), the total
    $O(K \log N)$ is much smaller than the $O(k \log N)$ baseline cost,
    not comparable to it. Concretely, this pass costs no more than
    $\frac{K}{k}$ of one baseline retrieval.
    \item \textbf{Alignment pass.} For each auditor $a_i$, the
    top-$k$ result against $\mathcal{M}_t$ is already produced by
    baseline retrieval (the auditor was about to use $q_i$ anyway),
    so calibration reuses it for free. Only the marginal effect of
    $\Delta_t$ needs to be computed: we score the $K$ new vectors
    against $q_i$ in $O(K)$ time and merge them into the cached
    top-$k$ in $O(k)$ time, then compute $d_{\text{RBO}}$ in $O(k)$.
    Total per auditor: $O(K + k)$, with no $\log N$ factor. Across
    auditors: $O(|\mathcal{A}_t|(K + k))$.
\end{itemize}
Total added cost: $O(K \log N + |\mathcal{A}_t|(K + k))$. Since both
$K$ and $|\mathcal{A}_t|(K + k)$ are independent of $N$ except for the
single $K \log N$ term (which itself is much smaller than baseline
because $K \ll k$ in practice), calibration overhead is sub-baseline.

\textbf{Graph case.}
Baseline graph retrieval traverses up to $L_{\max}$ hops at cost
$O(d^{L_{\max}})$, where $d$ is the average degree and $L_{\max} =
\lceil \log |\mathcal{V}_t| \rceil$. Calibration adds three pieces:
\begin{itemize}
    \item \textbf{Local check.} A neighborhood scan over $u$'s
    existing edges costs $O(d)$ per new edge, so $O(|\mathcal{E}_t^\Delta| \, d)$
    in total. Since baseline retrieval is $O(d^{L_{\max}})$ with
    $L_{\max} \geq 2$, this is smaller than baseline by a factor of
    $d^{L_{\max} - 1}$.
    \item \textbf{Path check.} Bidirectional BFS gives
    $O(d^{L_{\max}/2})$ per edge, which is the square root of the
    baseline's $O(d^{L_{\max}})$. Even summed over
    $|\mathcal{E}_t^\Delta|$ edges, this stays well below baseline as
    long as $|\mathcal{E}_t^\Delta| \ll d^{L_{\max}/2}$, true for any
    realistic incremental update.
    \item \textbf{Alignment pass.} A reachability check costs
    $O(\log |\mathcal{V}_t|)$ per auditor by reusing the adjacency
    index already loaded for baseline retrieval. Total:
    $O(|\mathcal{A}_t| \log |\mathcal{V}_t|)$.
\end{itemize}
All three terms are smaller than the baseline $O(d^{L_{\max}})$ by at
least one factor of $d^{L_{\max}/2}$, so calibration cost is again
sub-baseline.

\textbf{No extra LLM calls.}
The dominant cost in any agent-based system is the LLM forward pass.
\system makes zero additional LLM calls. Probes reuse queries and
walks the auditors generate during normal debate, and every scoring
step is algorithmic (ANN lookups, BFS, RBO). The asymptotic added
cost above is therefore the full added cost.

\textbf{Complexity table.}
Table~\ref{tab:complexity_comparison} states the asymptotic cost of
every stage and the ratio between added cost and the baseline
retrieval cost. The ratio is shown in the rightmost column and is
strictly less than $1$ across all stages under the regimes named in
the caption.

\begin{table}[h]
\caption{Time complexity of each calibration stage and the ratio of
added cost to baseline retrieval cost. The ratios are small under
typical operating regimes: $K \ll k$ (delta size much smaller than
top-$k$ width), $|\mathcal{A}_t|$ a small constant, $|\mathcal{V}_t|$
a small graph, $|\mathcal{E}_t^\Delta| \ll d^{L_{\max}/2}$
(incremental updates). Under these conditions every added stage is
sub-baseline, and total calibration overhead remains marginal
relative to one round of baseline retrieval.}
\centering
\small
\resizebox{\textwidth}{!}{
\begin{tabular}{lllll}
\toprule
\textbf{Arch.} & \textbf{Stage} & \textbf{Baseline cost} &
\textbf{\system added cost} & \textbf{Ratio (added / base)} \\
\midrule
\multirow{4}{*}{Embedding}
 & Retrieval / query & $O(k \log N)$ & --- & --- \\
 & Detection pass & --- & $O(K \log N)$ & $K / k \ll 1$ \\
 & Alignment pass & --- & $O(|\mathcal{A}_t|(K + k))$ &
   $|\mathcal{A}_t|(K{+}k) / (k \log N) \ll 1$ \\
 & Credibility / re-weight & --- & $O(|\mathcal{A}_t| + K)$ &
   negligible \\
\midrule
\multirow{4}{*}{Graph}
 & Retrieval / query & $O(d^{L_{\max}})$ & --- & --- \\
 & Local check & --- & $O(|\mathcal{E}_t^\Delta| \, d)$ &
   $|\mathcal{E}_t^\Delta| / d^{L_{\max} - 1} \ll 1$ \\
 & Path check & --- & $O(|\mathcal{E}_t^\Delta| \, d^{L_{\max}/2})$ &
   $|\mathcal{E}_t^\Delta| / d^{L_{\max}/2} \ll 1$ \\
 & Alignment & --- & $O(|\mathcal{A}_t| \log |\mathcal{V}_t|)$ &
   $|\mathcal{A}_t| \log |\mathcal{V}_t| / d^{L_{\max}} \ll 1$ \\
\bottomrule
\end{tabular}}
\label{tab:complexity_comparison}
\end{table}

Two observations close the analysis. \textbf{(i)} Every added stage
shares the same asymptotic terms as the baseline ($\log N$ for
embeddings, $d^{L_{\max}}$ or smaller for graphs), so calibration
scales with the same primitives the baseline uses. \textbf{(ii)} The
multiplicative factors in front of those terms ($K$, $|\mathcal{A}_t|$,
$|\mathcal{E}_t^\Delta|$) are all independent of memory size and stay
small in practice, which is what makes the added cost marginal rather
than merely comparable.

%-------------------------------------------------------------
\subsection{Memory and Token Overhead}
\label{app:memory_overhead}
%-------------------------------------------------------------

\textbf{Memory overhead.}
\system stores three additional pieces of state beyond the base
memory:
\begin{itemize}
    \item One scalar trust weight $\rho \in [0,1]$ per committed
    entry. For a memory of $N$ entries, this is $N$ floating-point
    numbers, or roughly $4N$ bytes at single precision. Compared to
    the embedding storage of $4dN$ bytes (for $d$-dimensional
    embeddings), this is a $1/d$ relative overhead, which is
    negligible at typical $d \in [768, 1536]$.
    \item One credibility weight $w_i^{(t)} \in [0,1]$ per auditor.
    For $|\mathcal{N}|$ agents, this is $|\mathcal{N}|$ scalars,
    constant in $N$.
    \item A bounded history buffer of recent calibration scores for
    computing $\rho^*$ and $\eta_t$. We maintain the last $W$ scores
    using a fixed-size ring buffer, so memory is $O(W)$ regardless of
    how many rounds have elapsed.
\end{itemize}
Total: an additional $O(N + W + |\mathcal{N}|)$ scalars on top of the
base memory, dominated by the trust-weight column.

\textbf{Token overhead.}
We emphasize again that \system makes zero additional LLM calls. The
contributor's $\Delta_t$ generation, the auditors' query and
path-walking activity, and the agent debate itself are all unchanged.
The probes \system uses are exactly the queries and walks the
auditors would issue without calibration. There is no calibration
prompt, no auditor-side reasoning solicited, and no contributor-side
justification generated. Token consumption per round is therefore
identical to an uncalibrated multi-agent debate of the same
configuration.

\textbf{Index overhead.}
For embedding memory we reuse the existing ANN index. Calibration
queries it but does not modify it beyond standard insertion, so no
new index structure or rebuild step is added. For graph memory we
reuse the existing adjacency structure plus a small side index
$\mathcal{R}_{\text{compat}}$ of relation co-occurrence counts. The
side index has size $O(|\mathcal{R}|^2)$ where $|\mathcal{R}|$ is the
number of relation types ($\leq 10^3$ even for large knowledge
graphs), so it adds at most a few megabytes regardless of
$|\mathcal{V}_t|$. Both index costs are constant in the memory size
and therefore do not enter the asymptotic analysis of
\S\ref{app:complexity}.

%-------------------------------------------------------------
\subsection{Compatibility Set Estimation}
\label{app:compat_estimation}
%-------------------------------------------------------------

The compatibility set $\mathcal{R}_{\text{compat}}(r)$ used in the
graph-memory local check is learned from the existing graph rather
than supplied as an external ontology. We compute it as follows.

For each pair of relation types $(r, r') \in \mathcal{R} \times
\mathcal{R}$, we count the number of nodes $u \in \mathcal{V}_t$
where $u$ has both an outgoing edge of type $r$ and an outgoing edge
of type $r'$. Call this count $c(r, r')$. We define the
co-occurrence score:
\begin{equation}
    \text{cooc}(r, r') = \frac{c(r, r')}{c(r) \cdot c(r') /
    |\mathcal{V}_t|}
\end{equation}
where $c(r) = \sum_{r'} c(r, r')$ is the marginal occurrence count
of relation $r$. This is the standard pointwise mutual information
form, normalized so that random co-occurrence yields
$\text{cooc} = 1$.

The compatibility set is then:
\begin{equation}
    \mathcal{R}_{\text{compat}}(r) = \{r' \in \mathcal{R} :
    \text{cooc}(r, r') \geq \tau_t\}
\end{equation}
where $\tau_t$ is the running median of all $\text{cooc}$ values in
the current graph (again adaptive, no manual setting). At cold start
(empty graph), $\mathcal{R}_{\text{compat}}(r) = \mathcal{R}$ for all
$r$, meaning the local check is uninformative until the graph has
enough edges to estimate co-occurrence. This degrades the detection
signal during warm-up but does not produce false rejections, because
$s_{\text{local}}$ stays at $1$ when the set is universal.

\textbf{Update cost.}
$\mathcal{R}_{\text{compat}}$ is recomputed incrementally. Each new
edge $(u, r, v)$ updates at most $|\mathcal{N}(u)|$ entries of $c(r,
\cdot)$, so per-edge update cost is $O(|\mathcal{N}(u)|)$. Summed
over the new-edge set, the total $O(|\mathcal{E}_t^\Delta| \cdot d)$
matches the local-check term in Table~\ref{tab:complexity_comparison}
and is therefore absorbed into the marginal-cost budget already
analyzed in \S\ref{app:complexity}. No separate latency cost is
incurred for maintaining $\mathcal{R}_{\text{compat}}$.

\textbf{Storage cost.}
$\mathcal{R}_{\text{compat}}$ is stored as a $|\mathcal{R}| \times
|\mathcal{R}|$ co-occurrence matrix, where $|\mathcal{R}|$ is the
number of relation types. Even at $|\mathcal{R}| = 10^3$ (large for a
knowledge graph), this is $10^6$ scalars, i.e., a few megabytes,
independent of $|\mathcal{V}_t|$. Compared to the base graph
storage of $O(|\mathcal{V}_t| + |\mathcal{E}_t|)$ entries, this is a
constant additive overhead and does not affect any complexity
bound.

%-------------------------------------------------------------
\subsection{Bidirectional BFS Depth and Cost}
\label{app:graph_bfs}
%-------------------------------------------------------------

We chose $L_{\max} = \lceil \log |\mathcal{V}_t| \rceil$ for the
bidirectional BFS in the path check. This choice rests on two
empirical regularities of knowledge graphs.

First, real-world knowledge graphs satisfy a small-world property:
the typical shortest-path distance between two random nodes scales
as $O(\log |\mathcal{V}|)$. This holds for ConceptNet, Wikidata,
Freebase, and the graph memories used by \cite{anokhin2024arigraph}
and \cite{zhang2025g}. Setting $L_{\max}$ to this scale captures
essentially all paths the contributor could plausibly justify, while
truncating exponential blowup of BFS at deeper levels.

Second, bidirectional BFS searches from both endpoints simultaneously
and meets in the middle, so the effective branching factor is
$d^{L_{\max}/2}$ rather than $d^{L_{\max}}$. With $L_{\max} = \log
|\mathcal{V}_t|$ this gives a worst-case cost of $\sqrt{|\mathcal{V}_t|}
\cdot \text{poly}(d)$, which is sub-linear in graph size.

\textbf{Failure modes.}
If the graph is unusually dense or has a large diameter, $L_{\max} =
\log |\mathcal{V}_t|$ may either over-explore or miss valid paths. We
detect this online by tracking the running ratio
$|\mathcal{P}_{uv}| / \text{expected paths}$ across rounds. If the
ratio drifts outside $[0.5, 2.0]$, we adjust $L_{\max}$ by $\pm 1$.
The adjustment changes the BFS cost by at most a factor of $d^{1/2}$
(the same bidirectional speedup applies), so the marginal-cost
analysis in Table~\ref{tab:complexity_comparison} still holds within
a constant factor.

%-------------------------------------------------------------
\subsection{Retrieval Complexity Under Trust Weighting}
\label{app:retrieval_complexity}
%-------------------------------------------------------------

We claimed in \S\ref{ssec:retrieval} that trust weighting adds zero
asymptotic cost to retrieval (not just marginal cost, but no added
cost at all). We justify this for both architectures.

\textbf{Embedding retrieval.}
The trust-weighted vector $\tilde{\mathbf{e}}_k = \sqrt{\rho_k} \cdot
\mathbf{e}_k$ is computed once at index-load time and stored in place
of $\mathbf{e}_k$. Top-$k$ retrieval against
$\{\tilde{\mathbf{e}}_k\}$ uses the same ANN index structure as
retrieval against $\{\mathbf{e}_k\}$, with the same $O(k \log N)$
cost per query. When an entry's $\rho_k$ is updated, we re-scale that
single vector in $O(d)$ and re-insert into the ANN index in $O(\log
N)$, identical to a standard memory write. No new cost term enters
Table~\ref{tab:complexity_comparison}.

\textbf{Graph retrieval.}
For path queries, the trust-weighted path strength $\prod_{e \in \pi}
\rho_e$ is computed lazily during traversal: each edge access
multiplies the running product by $\rho_e$. This is one extra
multiplication per visited edge, an $O(1)$ overhead on top of the
$O(1)$ baseline edge-access cost. Total query complexity stays at
$O(d^{L_{\max}})$, identical to baseline. When an edge's $\rho_e$ is
updated, we modify a single entry in the edge attribute table in
$O(1)$.

%-------------------------------------------------------------
\subsection{Implementation Notes}
\label{app:impl_notes}
%-------------------------------------------------------------

We provide a few practical notes for re-implementation.

\textbf{ANN index choice.}
We use HNSW \citep{malkov2018efficient} with $M = 16$ and $\text{ef} =
200$ for the embedding case. The detection pass requires range
queries (count neighbors within a similarity radius), which HNSW does
not natively support, so we issue a top-$k$ query with a generous
upper bound (e.g., $k_{\text{filter}} = 100$) and filter by the
adaptive radius post-hoc. The post-hoc filter is $O(k_{\text{filter}})$
per call and does not change the $O(\log N)$ asymptotic of the
underlying ANN query, so it is absorbed into the $O(K \log N)$
detection-pass cost reported in
Table~\ref{tab:complexity_comparison}.

\textbf{RBO truncation.}
Rank-biased overlap is defined as an infinite series, but it
converges geometrically. We truncate at depth $k$ (the same $k$ used
for top-$k$ retrieval) with persistence parameter $p_{\text{RBO}} =
0.9$, the default in the original RBO paper. This is a property of
RBO itself and not a hyperparameter introduced by \system.

\textbf{Probe sampling.}
For embedding memory, $Q_i$ is the auditor's pending-query queue
(typically size $1$ to $5$ per round). For graph memory, $\pi_i$ is
the auditor's most recent traversal path. In both cases, we sample
uniformly using the system's standard random number generator with a
seed reset at each round to prevent adaptive auditors from predicting
future samples from past behavior.

\textbf{Numerical stability.}
We clamp all $\rho$, $w_i$, $\rho_{\text{geo}}$, and $\rho_{\text{ret}}$
values into $[10^{-6}, 1 - 10^{-6}]$ before taking products or square
roots, to avoid underflow when scores accumulate over many rounds.
This is standard numerical hygiene and not a tunable parameter.